\documentclass[../main.tex]{subfiles}

\begin{document}

\section{Exersice}


\begin{exercise}{Calculate}{}
Let $V$ be a vector space of finite dimension with basis $(e_1, \dots, e_n)$ and let $V^*$ denotes the dual space with dual basis $(e_1^*, \dots, e_n^*)$. Please prove in two different ways that for every $(\alpha, \beta) \in \Lambda^p V^* \times \Lambda^q V^*$,
\[
\alpha \wedge \beta = (-1)^{pq} \beta \wedge \alpha.
\]
One way is to use definition, the other is to calculate in basis using the next exercise.
\end{exercise}

\begin{proof}
   \begin{enumerate}
    \item   \[
     \alpha \wedge \beta = \frac{\left( p+ q \right)!  }{ p!q!} \mathrm{Alt}\left( \alpha \otimes \beta  \right) 
     \] \[
     \beta \wedge \alpha =  \frac{\left( p+ q \right)!  }{p!q! } \mathrm{Alt}\left( \beta \otimes \alpha  \right)  
     \]Only need to show that \[
     \mathrm{Alt}\left( \alpha \otimes \beta  \right)= \left( -1 \right)^{pq} \left( \beta \otimes \alpha  \right)   
     \] Where \[
     \mathrm{Alt}\left( \alpha \otimes \beta  \right)= \frac{1 }{\left( p+ q \right)!  }\sum _{ \sigma \in  S_{p+ q}} \left( \mathrm{sgn} \sigma  \right)    \sigma \left(  \alpha \otimes \beta  \right) 
     \]Let \(  \tau = \left( p+ 1,\cdots ,p+ q, 1,2,\cdots ,p \right)   \) , then \(  \mathrm{sgn}\left( \tau  \right)= \left( -1 \right)^{pq}    \)   . We have \[
    \begin{aligned}
     \mathrm{Alt}\left( \alpha \otimes \beta  \right)&=  \frac{1 }{\left( p+ q \right)!  } \sum _{ \sigma \in S_{p+ q}}\left( \mathrm{sgn} \sigma  \right)  \sigma\circ  \tau \left( \beta \otimes \alpha  \right)     \\ 
      &= \left( -1 \right)^{pq}\frac{1 }{\left( p+ q \right)!  } \sum _{ \sigma \in S_{p+ q}} \left( \mathrm{sgn} \sigma  \right)\left( \mathrm{sgn}\tau  \right) \left(  \sigma \circ \tau  \right)\left( \beta \otimes \alpha  \right)\\ 
       &= \left( -1 \right)^{pq}\frac{1 }{\left( p+ q \right)!  }  \sum _{  \sigma ^{\prime} \in S_{p+ q}} \left( \mathrm{sgn} \sigma ^{\prime}  \right)  \sigma ^{\prime} \left( \beta \otimes \alpha  \right)\\ 
        &= \left( -1 \right)^{pq} \mathrm{Alt}\left( \beta \otimes \alpha  \right)            
    \end{aligned}
     \]
     \item Suppose \[
     \alpha =  \sum _{i_1 \cdots i_{p}}a_{i_1\cdots i_{p}}e_{i_1}^{*}\wedge \cdots \wedge e_{i_{p}}^{*}
     \]  \[
     \beta = \sum _{j_1\cdots j_{q}} b_{j_1\cdots j_{q}} e_{j_1}^{*}\wedge \cdots \wedge e_{j_{q}}^{*}
     \] Then \[
     \alpha \wedge \beta =  \sum _{i_1\cdots i_{p}} a_{i_1\cdots i_{p}}b_{j_1\cdots j_{q}}\left( e^{*} \right)_{I} \wedge \left( e^{*} \right)_{J}= \sum _{I,J} a_{I}b_{J} e_{IJ}  
     \] \[
     \beta \wedge {\alpha }= \sum _{I,J} a_{I}b_{J}e_{JI}= \left( -1 \right)^{pq} \sum _{I,J}a_{I}b_{J}e_{IJ} = \left( -1 \right)^{pq}\alpha \wedge \beta  
     \]
   \end{enumerate}
   
\end{proof}
    

\begin{exercise}{}{}
Let $V$ be a vector space of finite dimension and let $V^*$ denotes the dual space. Prove that if $(\theta_1, \dots, \theta_r)$ is a family of r covectors in $V^*$ which are \dfntxt{linearly dependent}, then the following wedge product:
\[
\theta_1 \wedge \dots \wedge \theta_r = 0.
\]
\end{exercise}


\begin{proof}
    
    Since \(   \theta _1 \wedge \cdots  \theta _{r}  \) is multi-linear and alternative,  take any vector \(   v^1,\cdots,v^r   \) , we have \[
    \begin{aligned}
    \theta _1 \wedge \cdots \wedge  \theta _{r}\left(  v^1,\cdots,v^r  \right)&= \det \left(  \theta _{i}\left( v^{j} \right)  \right) _{ij}
    \end{aligned} 
    \]  Since \(   \theta _1 ,\cdots , \theta _{r}  \) are linealy dependent , the  row vectors of the matrix  \(  \left(  \theta _{i}\left( v^{j} \right)  \right)_{ij}   \) are as well. Thus \(  \det \left(  \theta _{i}\left( v^{j} \right)  \right)_{ij}= 0   \) , and \(   \theta _1 \wedge \cdots  \theta _{r}\left(  v^1,\cdots,v^r  \right)= 0   \).    Since \(   v^1,\cdots,v^r   \) are arbitrary ,  \(   \theta _1 \wedge \cdots \wedge  \theta _{r}= 0  \).   
\end{proof}

\begin{exercise}{Basis of Grassmann algebra}{}
Let $V$ be a vector space of finite dimension with basis $(e_1, \dots, e_n)$ and let $V^*$ denotes the dual space with dual basis $(e_1^*, \dots, e_n^*)$. Please prove that
\[ (e_{i_1}^* \wedge \dots \wedge e_{i_k}^*)_{1 \le i_1 < \dots < i_k \le n} \]
forms a basis of $\Lambda^k V^*$.
\end{exercise}

\begin{proof}
   For \(  \alpha  \in  \Lambda ^{k}V^{*}  \), and \(  I= \left(  i_1,\cdots,i_k  \right)   \) . Define \[
    \alpha _{I}: = \alpha \left( e_{i_1},\cdots ,e_{i_{k}} \right) 
    \]  If  \(  I   \) has a reputed index , then \(  \alpha _{I}= 0  \)   , and if \(  J=  \sigma I \) ,  \(  \alpha _{J}= \left( \mathrm{sgn} \sigma  \right)\alpha _{I}   \) .  We have \[
   \sum _{I}^{^{\prime} }\alpha _{I} e^{*}_{I}\left( e_{j_1},\cdots ,e_{j_{k}} \right)= \sum _{I}^{^{\prime} }\alpha _{I}   \delta _{J}^{I}= \alpha _{J}= \alpha \left( e_{j_1},\cdots ,e_{j_{k}} \right) 
    \]Thus  \(  \alpha = \sum _{I}^{^{\prime} }\alpha _{I} e_{I}^{*}  \)  . Which shows that  \(  e^{*}_{i_1},\cdots ,e^{*}_{i_{k}}, 1\le i_1< \cdots < i_{k}\le n  \) spands the  \(   \Lambda ^{k}V^{*}  \).  To show that they are linealy independent.   Supppose that \[
    \sum _{I} k_{I}e ^{*}_{I}= 0
    \] Both sides act on \(  e_{J}   \), we have \(  k_{J}= 0  \).   
\end{proof}

\begin{exercise}{Grassmann algebra versus Polynomial functions on V.}{}
Let $V$ be a vector space of finite dimension and let $V^*$ denotes the dual space.
A p-multilinear map $\alpha$ on $V$ is \dfntxt{symmetric} if
\[ \alpha(v_1, \dots, v_n) = \alpha(v_{\sigma(1)}, \dots, v_{\sigma(n)}) \]
for all permutations $\sigma$.
Such $\alpha$ defines a homogeneous polynomial $\tilde{\alpha}$ defined by
\[ \tilde{\alpha}(v) = \alpha(v, \dots, v) \]
for all $v \in V$.
Then show that for symmetric multilinear maps $\alpha, \beta$ on $V$ of respective p, q, the product $\alpha \circ \beta$ defined by
\[ \alpha \circ \beta(v_1, \dots, v_{p+q}) := \frac{1}{p!q!} \sum_{\sigma \in S_{p+q}} \alpha(v_{\sigma(1)}, \dots, v_{\sigma(p)}) \beta(v_{\sigma(p+1)}, \dots, v_{\sigma(p+q)}) \]
is valued into symmetric multilinear map on $V$ of degree $p+q$ and that
\[ \alpha \circ \beta(v, \dots, v) = \alpha(v, \dots, v) \beta(v, \dots, v). \]
Could you compare with $\Lambda^\bullet V^*, \wedge$ ?
\end{exercise}

\begin{exercise}{Basic important examples of smooth manifolds}{}
Please prove that the following topological spaces have a smooth manifold structure :
\begin{itemize}
    \item The Euclidean sphere $S^n \subset \mathbb{R}^{n+1}$.
    \item The matrix groups $GL_n(\mathbb{R}), SL_n(\mathbb{R}), O_n(\mathbb{R})$, these are Lie groups of matrice. Each time could you tell me what is the dimension of these groups.
\end{itemize}
\end{exercise}
\begin{proof}
 \begin{itemize}
    \item Define \(  F  \): \[
    F\left( x^{1},\cdots ,x^{n+ 1} \right)= x_1^{2}+ \cdots + x_{n+ 1}^{2} 
    \]  Then \[
    S^{n}=  F^{-1} \left( 1 \right) 
    \] \[
    \,\mathrm{d} F_{x}= \nabla F\left( x \right)= \left( 2x_1,\cdots ,2_{x_{n+ 1}} \right)    \neq 0 ,\forall x \in \mathbb{R} ^{n+ 1}
    \]By regular value theorem, \(  S^{n}\subseteq \mathbb{R} ^{n+ 1}  \)  is a regular embeded submanifold  of \(  \mathbb{R} ^{n+ 1}  \) , which  itself is a  manifold.  
    
\item 
  \begin{enumerate}
    \item \(  \mathrm{Mat}\left( n,\mathbb{R}  \right)\simeq \mathbb{R} ^{2}   \)是\(  n^{2}  \)维流形.  
    \item \(  \operatorname{GL} \left( n,\mathbb{R}  \right)= \left\{ M\in \mathrm{Mat}\left( n,\mathbb{R}  \right):\det M\neq 0  \right\}   \)作为 \(  \mathrm{Mat}\left( n,\mathbb{R}  \right)   \)的一个开子集也是一个\(  n^{2}  \)维流形.    
    \item  \(  \mathrm{SL}\left( n,\mathbb{R}  \right)= \left( M\in \mathrm{Mat}\left( n,\mathbb{R}  \right)  \right):\det M= 1    \)   \[
    \begin{aligned}
    \det : \mathrm{Mat}\left( n,\mathbb{R}  \right)&\to \mathbb{R} \\ 
     M&\mapsto \det M  
    \end{aligned}
    \]  \(  SL  \left( n,\mathbb{R}  \right)= \det ^{-1} \left( I \right)  \). 取一个 \(  M \in \mathrm{Mat}\left( n,\mathbb{R}  \right)   \), \(  \det M= 1  \). 要证明 \(  \left( \,\mathrm{d} \left( \det  \right)  \right)_{M}: \mathbb{R} ^{n^{2}}\to \mathbb{R}    \)   . 按分量求偏导,得到 \[
    \frac{\partial }{\partial a_{ij}} \left( \det M \right)=  A_{ij}(\text{代数余子式})  
    \]由于 \(  \det M= 1  \).  至少有一个 \(  A_{ij}\neq 0  \). 从而  \(  \left( \,\mathrm{d} \left( \det  \right)  \right)_{M}   \)是满的. 
  \item 定义 \(  P: \mathrm{Mat}\left( n,\mathbb{R}  \right)\to  \mathrm{SMat}\left( n,\mathbb{R}  \right)(\text{对称矩阵}) \simeq \mathbb{R} ^{\frac{n\left( n+ 1 \right)  }{2 } }   \) .  \[
  P\left( A \right): = A^{\top}A 
  \]则 \(  O\left( n,\mathbb{R}  \right)= P^{-1} \left( I \right)   \) . 对于 \(  A \in O\left( n,\mathbb{R}  \right)   \). \[
  \left( \,\mathrm{d} P \right)_{A}: \mathrm{Mat}\left( n,\mathbb{R}  \right)\to \mathrm{SMat}\left( n,\mathbb{R}  \right)   
  \]与上例不同, 这里用分量的形式去计算 Jacobi是看不清楚的, 不过由于 \(  P  \)本身足够简单, 我们可以用微分最原始的含义来观察. 考虑  \[
  P\left( A+  \varepsilon B \right)-P\left( A \right)=  \varepsilon \left( \,\mathrm{d} P \right)_{A}\left( B \right)+ O\left(  \cdots   \right)     
  \]带入得到 \[
  \left( A^{T}+  \varepsilon B^{T} \right)\left( A+  \varepsilon B \right)-A^{\top}A=    \varepsilon \left(B^{\top}A+ A^{\top}B \right)+ O\left(  \varepsilon  \right)  
  \]于是 \(  \left( \,\mathrm{d} P \right)_{A}   \left( B \right)= B ^{\top}A+ A^{\top}B\) . 只需要证明这个关于 \(  B  \)的映射是满的. 为此, 解方程 \[
  A^{\top}B+ B^{\top}A= C
  \] 取 \(  B = \frac{1}{2}AC  \), 利用 \(  A,C  \)对称,   带入得到 \[
  \frac{1}{2}A^{\top}AC+ \frac{1}{2}C^{\top}A^{\top}A= \frac{1}{2}\left( C+ C^{\top} \right)= C 
  \]  
  故 \(  O\left( n,\mathbb{R}  \right)   \)是 \(  \mathrm{Mat}\left( n,\mathbb{R}  \right)   \)的一个 \(  \frac{n\left( n-1 \right)  }{2 }   \)维正则子流形. 注意到对于 \(  A \in O\left( n,\mathbb{R}  \right)   \) , \(  \det A= 1  \)或 \(  \det A= -1  \) , 可以窥见事实上 \(  O\left( n,\mathbb{R}  \right)   \)有两个连通分支.  
  \end{enumerate}
  
 \end{itemize}
  
\end{proof}
\begin{exercise}{Quotient manifolds}{}
Prove that the following objects are quotient manifolds
\begin{itemize}
    \item $\mathbb{R}P^n$ obtained by quotienting $S^n$ by antipodal map, show that it is nonorientable by a picture. Could you explain by a picture that the neighborhood of the line at infinity is homeomorphic to a Möbius ribbon ?
    \item Prove that the torus $\mathbb{T}^d = \mathbb{R}^d/\mathbb{Z}^d$ is a smooth manifold.
\end{itemize}
\end{exercise}
\begin{proof}

    
    \begin{itemize}


        \item 
        \begin{enumerate}
            \item See \(  \mathbb{R} P^{n}  \) is \(  S^{2}  /\left( x\sim -x \right) \).  The orentation changes at the equator.    The map \(  A: S^{2}\to S^{2} . A\left( x \right)= -x   \) induces an idendity map on \(  \mathbb{R} P^{2}  \), \(  \,\mathrm{d} A  \) send an orentation basis \(  \left\{ e_1,e_2 \right\}  \) to one opposite to it .  
            \item Use $\mathbb{R}P^2 = D^2 / \{x \sim -x \text{ for } x \in \partial D^2 = S^1\}$. Then the neighborhood of the line at infinity is a thin ring near \(   \partial D^{2}  \) with some quotient relation. The quotient is just the same at that of Möbius .
        \end{enumerate}
        
    \end{itemize}
    
\end{proof}
\begin{exercise}{Vector fields are derivation}{}
Let M be a $C^\infty$ manifold. Show that the left $C^\infty$ module of derivation of the ring $C^\infty(M)$ i.e. linear maps $D: C^\infty(M) \to C^\infty(M)$ satisfying
\[ D(fg) = (Df)g + f(Dg) \]
can be identified with vector fields on M, smooth sections of TM.
\end{exercise}
\begin{proof}
    
    Define \[
    D\left( x \right)f=  \left( Df \right)\left( x \right)        
    \] Then \[
    D\left( x \right)\left( fg \right)=   
    \]
\end{proof}
\end{document}